% 绝对连续的多区间控制：F(x)=sqrt(x) 于 [0,4]；四个不交区间 (0.2,0.5),(1.4,1.8),(2.6,2.9),(3.4,3.6)，总长 1.2；对应增量 0.260,0.158,0.090,0.053，总和约 0.562。纵轴放大 1.2 倍绘制。
\begin{tikzpicture}[>=Stealth,
  ax/.style={draw=FigureStroke,line width=.6pt,->},
  guide/.style={draw=FigureMuted,dashed,line width=.4pt},
  note/.style={font=\FigureNoteFont\fontsize{7.5}{9}\selectfont,text=FigureMuted}]
\begin{scope}[xscale=1.6,yscale=1.2]
\draw[ax] (-.1,0)--(4.35,0) node[right] {$x$};
\draw[ax] (0,-.05)--(0,2.25) node[above] {$F(x)=\sqrt x$};
\draw[FigureBlue,line width=1.1pt,domain=0:4,samples=120] plot (\x,{sqrt(\x)});
\foreach \a/\b in {0.2/0.5,1.4/1.8,2.6/2.9,3.4/3.6}{
  \draw[FigureGreen,line width=2.4pt] (\a,0)--(\b,0);
  \draw[FigureRed,line width=2.4pt] (0,{sqrt(\a)})--(0,{sqrt(\b)});
  \draw[guide] (\a,0)--(\a,{sqrt(\a)})--(0,{sqrt(\a)});
  \draw[guide] (\b,0)--(\b,{sqrt(\b)})--(0,{sqrt(\b)});}
\foreach \x in {1,2,3,4}{\node[below,font=\fontsize{7.5}{9}\selectfont] at (\x,-.05) {$\x$};}
\end{scope}
\end{tikzpicture}
